% Part: many-valued-logic
% Chapter: syntax-and-semantics
% Section: sublogics

\documentclass[../../../include/open-logic-section]{subfiles}

\begin{document}

\olfileid{mvl}{syn}{sub}

\olsection{$\LogCL$కు ఉపతర్కాలుగా బహుమూల్య తర్కాలు}

సాధారణంగా వాడే బహుమూల్య తర్కాల మాత్రికల్లో, $\{\True, \False\}$లోని
విలువలను వాదాలుగా ఇచ్చినప్పుడు సత్యమూల్య ప్రమేయాలు సాంప్రదాయిక ప్రమేయాల
ఫలితాలనే ఇస్తాయి. అంటే, ఈ తర్కాల్లో
$x \in \{\True, \False\}$ అయితే $\tf{\lnot}[\Log L](x) =
\tf{\lnot}[\LogCL](x)$; అలాగే $\star$ అనేది $\land$, $\lor$,
$\lif$లలో ఏదైనా, $x, y \in \{\True, \False\}$ అయితే
$\tf{\star}[\Log L](x,y) = \tf{\star}[\LogCL](x,y)$.
అంటే $\lnot$, $\land$, $\lor$, $\lif$లకు చెందిన
సత్యమూల్య ప్రమేయాలను $\{\True,\False\}$కు పరిమితం చేస్తే
సాంప్రదాయిక సత్యమూల్య ప్రమేయాలే వస్తాయి.

\begin{prop}\ollabel{prop:mvl-cl}
  బహుమూల్య తర్కం~$\Log L$ భాషలో $\lnot$, $\land$, $\lor$,
  $\lif$ సంయోజకాలు ఉన్నాయనీ, $\True, \False \in
  V$ అనీ అనుకుందాం. దాని సత్యమూల్య ప్రమేయాలు కింది షరతులను
  సంతృప్తిపరుస్తాయని కూడా అనుకుందాం:
  \begin{enumerate}
    \item\ollabel{prop:not} $\tf{\lnot}[\Log L](x) = \tf{\lnot}[\LogCL](x)$, $x =
    \True$ లేదా $x = \False$ అయితే;
    \item\ollabel{prop:land} $\tf{\land}[\Log L](x,y) = \tf{\land}[\LogCL](x,y)$,
    \item\ollabel{prop:lor} $\tf{\lor}[\Log L](x,y) = \tf{\lor}[\LogCL](x,y)$,
    \item\ollabel{prop:lif} $\tf{\lif}[\Log L](x,y) = \tf{\lif}[\LogCL](x,y)$.
  \end{enumerate}
  చివరి మూడు సమీకరణాలూ $x, y \in \{\True, \False\}$ అయితే వర్తిస్తాయి.
  అప్పుడు, ప్రతి ప్రతిజ్ఞావాక్య చరానికీ $\pAssign
  v(p) \in \{\True,\False\}$ అయ్యేలా~$V$లోని విలువలను ఇచ్చే
ఏ సత్యమూల్య కేటాయింపు~$\pAssign v$నైనా తీసుకుందాం. ఈ నాలుగు
  సంయోజకాలతో మాత్రమే ప్రతిజ్ఞావాక్య చరాల నుంచి నిర్మించిన
  ప్రతి సూత్రానికీ $\pValue v[\Log L](!A) = \pValue
  v[\LogCL](!A)$.
\end{prop}

\begin{proof}
  $!A$ నిర్మాణంపై ఆగమనం ద్వారా నిరూపిద్దాం.
  \begin{enumerate}
  \item $!A \ident p$ పరమాణువైతే, $\pValue v[\Log L](!A) =
  \pAssign v(p) = \pValue
  v[\LogCL](!A)$.
  \item $!A \ident \lnot !B$ అయితే,
  \begin{align*}
    \pValue v[\Log L](!A) & = \tf{\lnot}[\Log L](\pValue v[\Log L](!B)) & &\text{\olref[val]{defn:pValue} ప్రకారం}\\
    & = \tf{\lnot}[\Log L](\pValue v[\LogCL](!B)) && \text{ఆగమన పరికల్పన ప్రకారం}\\
    & = \tf{\lnot}[\LogCL](\pValue v[\LogCL](!B))
&&\text{\olref{prop:not} పరికల్పన ప్రకారం,}\\
&&&\text{ఎందుకంటే $\pValue v[\LogCL](!B) \in \{\True, \False\}$,}\\
& = \pValue v[\LogCL](!A) &&\text{\olref[val]{defn:pValue} ప్రకారం}.
\end{align*}
\item $!A \ident (!B \land !C)$ అయితే,
\begin{align*}
  \pValue v[\Log L](!A) & = \tf{\land}[\Log L](\pValue v[\Log L](!B), \pValue v[\Log L](!C)) & &\text{\olref[val]{defn:pValue} ప్రకారం}\\
  & = \tf{\land}[\Log L](\pValue v[\LogCL](!B),\pValue v[\LogCL](!C)) && \text{ఆగమన పరికల్పన ప్రకారం}\\
  & = \tf{\land}[\LogCL](\pValue v[\LogCL](!B),\pValue v[\LogCL](!C))
&&\text{\olref{prop:land} పరికల్పన ప్రకారం,}\\
&&&\text{ఎందుకంటే $\pValue v[\LogCL](!B),\pValue v[\LogCL](!C) \in \{\True, \False\}$,}\\
& = \pValue v[\LogCL](!A) &&\text{\olref[val]{defn:pValue} ప్రకారం}.
\end{align*}
\end{enumerate}
$!A \ident (!B \lor !C)$, $!A \ident (!B \lif !C)$ సందర్భాలు
కూడా ఇదే విధంగా నిరూపించవచ్చు.
\end{proof}

\begin{cor}
  ఒక బహుమూల్య తర్కం \olref{prop:mvl-cl}లోని షరతులను
  సంతృప్తిపరచి, $\True \in V^+$, $\False \notin V^+$
  అయితే, పై నాలుగు సంయోజకాలతో చరాల నుంచి నిర్మించిన
  సాధారణ భాగానికి ${\Entails[\Log L]} \subseteq {\Entails[\LogCL]}$;
  అంటే, ఆ భాగంలోని సూత్రాలకు $\Gamma
  \Entails[\Log L] !B$ నుంచి $\Gamma \Entails[\LogCL] !B$ వస్తుంది.
  ముఖ్యంగా, ఆ భాగంలో $\Log L$కు సర్వసత్యమైన ప్రతి సూత్రమూ
సాంప్రదాయిక తర్కంలోనూ సర్వసత్యమే.
\end{cor}

\begin{proof}
  ప్రతివిపర్యయాన్ని నిరూపిస్తాం. $\Gamma \Entails/[\LogCL] !B$
  అనుకుందాం. అప్పుడు ఏదో \tetoken{సత్యమూల్య కేటాయింపు}{!!{valuation}}~$\pAssign v\colon \PVar \to
  \{\True, \False\}$ ఉండి, అన్ని $!A \in \Gamma$లకు
  $\pValue v[\LogCL](!A) = \True$, అలాగే
  $\pValue v[\LogCL](!B) = \False$. $\True,
  \False \in V$ కనుక \tetoken{సత్యమూల్య కేటాయింపు}{!!{valuation}}~$\pAssign v$
  తర్కం~$\Log L$కు కూడా \tetoken{ఒక సత్యమూల్య కేటాయింపు}{!!a{valuation}}.
  \olref{prop:mvl-cl} ప్రకారం, అన్ని $!A \in \Gamma$లకు
  $\pValue v[\Log L](!A) = \True$, ఇంకా
  $\pValue v[\Log L](!B) = \False$.
  $\True \in V^+$, $\False \notin V^+$ కనుక
  $\pSat{v}{\Gamma}[\Log L]$, $\pSat/{v}{!B}[\Log L]$;
  అంటే $\Gamma \Entails/[\Log L] !B$.
\end{proof}
\end{document}
